\chapter{Multimodale KI -- Vision, Audio und Video}
\label{chap:multimodal}

% ============================================================
% LERNZIELE
% ============================================================
\begin{lernziele}
Nach diesem Kapitel kannst du:
\begin{itemize}
    \item Multimodale Modelle (GPT-4o, Claude 4 Vision, Gemini 3) in Produktion einsetzen
    \item Bilder, Audio und Video als LLM-Input verarbeiten und strukturieren
    \item Eine Vision-RAG-Pipeline aufbauen (Bild-zu-Text + Vector Search)
    \item Multimodale Sicherheitsrisiken erkennen (Prompt Injection via Bild, OCR-Exploits)
    \item Die Tradeoffs zwischen Text-, Vision- und Audio-Modalitäten bewerten
\end{itemize}
\end{lernziele}

% ============================================================
% MOTIVATION
% ============================================================
\section{Motivation}

LLMs waren reine Textmaschinen. \textbf{GPT-4o, Claude 4 Vision und Gemini 3} verarbeiten Bilder, Audio und Video nativ. Mehr als 60\,\% der neuen Enterprise-KI-Anwendungen nutzen multimodale Eingaben (Gartner). Ein Support-Bot, der Screenshots liest, ein Code-Review-Tool, das Architekturdiagramme versteht, oder ein Qualitätsinspektor, der Kamerabilder auswertet -- das sind keine Zukunftsszenarien mehr, sondern aktuelle Produktionssysteme.

Der Shift von Text-zu-Multimodal ist der grösste Infrastrukturwandel im AI Engineering seit der Einführung von RAG. Er betrifft nicht nur die Modell-Wahl, sondern die gesamte Pipeline: Datenvorbereitung, Embedding, Retrieval, Sicherheit und Evaluation.

% ============================================================
% GRUNDLAGEN
% ============================================================
\section{Grundlagen multimodaler Modelle}

\subsection{Modalitäten und ihre Verarbeitung}

Ein multimodales Modell unterscheidet sich fundamental von einem Text-LLM: Es besitzt separate Encoder für jede Modalität (Vision, Audio, Text), die in einen gemeinsamen Representierungsraum projiziert werden. Die Entscheidungsebene (Transformer-Decoder) operiert auf dieser fusionierten Repräsentation.

\begin{verbatim}
   +--------+    +-----------+    +---------+
   | Text   |--->| Tokenizer |--->|         |
   +--------+    +-----------+    |         |
   +--------+    +-----------+    | Fusion  |    +----------+
   | Image  |--->| Vision    |--->| Encoder |--->| Decoder  |
   +--------+    | Encoder   |    |         |    | (LLM)    |
   +--------+    +-----------+    |         |    +----------+
   | Audio  |--->| Audio     |--->|         |
   +--------+    | Encoder   |    +---------+
                  +-----------+
\end{verbatim}

\begin{table}[H]
\centering
\begin{tabularx}{\linewidth}{lXcc}
\toprule
\textbf{Modell} & \textbf{Modalitäten} & \textbf{Vision Auflösung} & \textbf{Preis (Input)} \\
\midrule
GPT-4o & Text, Bild, Audio & 512$\times$512 -- 4096$\times$4096 & 2,50 \$/M Tokens \\
Claude 4 Vision & Text, Bild & bis 8192$\times$8192 & 3,00 \$/M Tokens \\
Gemini 3 Pro & Text, Bild, Audio, Video & nat. Auflösung & 1,25 \$/M Tokens \\
Llama 4 90B & Text, Bild & 1024$\times$1024 & Open Source \\
\bottomrule
\end{tabularx}
\caption{Multimodale Modelle im Vergleich}
\label{tab:multimodal-models}
\end{table}

\subsection{Token-Äquivalenz von Bildern}

Ein zentrales Konzept für das Cost-Engineering multimodaler Systeme: \textbf{Bilder werden in Tokens umgerechnet}. Ein 1024$\times$1024-Bild kostet bei GPT-4o ca. 2.000 Tokens -- unabhängig vom Inhalt. Ein leeres Bild kostet genauso viel wie ein komplexes Diagramm.

\begin{lstlisting}[language=Python, caption=Token-Kosten eines Bildes berechnen]
def image_token_cost(image_path: str, model: str = "gpt-4o") -> dict:
    """Berechnet die Token-Kosten eines Bildes fuer ein multimodales Modell."""
    from PIL import Image

    img = Image.open(image_path)
    w, h = img.size

    # GPT-4o: 170 Tokens pro 512x512 Tile, aufgerundet
    tiles_x = (w + 511) // 512
    tiles_y = (h + 511) // 512
    tiles = tiles_x * tiles_y

    base_tokens = 85    # Basis-Overhead pro Bild
    tile_tokens = 170   # Pro 512x512 Tile

    total_tokens = base_tokens + tiles * tile_tokens

    return {
        "dimensions": f"{w}x{h}",
        "tiles": tiles,
        "tokens": total_tokens,
        "estimated_cost_usd": total_tokens * 2.5 / 1_000_000,
    }

# Beispiel: Full-HD Screenshot (1920x1080)
cost = image_token_cost("screenshot.png")
print(f"{cost['tiles']} Tiles, {cost['tokens']} Tokens, "
      f"${cost['estimated_cost_usd']:.5f}")
# Ausgabe: 16 Tiles, 2805 Tokens, $0.007
\end{lstlisting}

\begin{praxishinweis}
Bild-Tokens summieren sich schneller als gedacht. Ein einzelnes 4K-Bild kostet bei GPT-4o high-detail rund 7.000 Tokens -- mehr als dieser gesamte Absatz. Vorverarbeitung lohnt: zuschneiden, skalieren, Rauschen reduzieren.
\end{praxishinweis}

% ===========================================================%
% VISION
% ============================================================
\section{Vision -- Bilder und Dokumente}

\subsection{Image-Input-Pipeline}

Die Verarbeitung von Bildern in einer LLM-Pipeline folgt einem standardisierten Muster: Bild laden, ggf. vorverarbeiten, Base64-kodieren oder per URL referenzieren, multimodalem Prompt hinzufügen.

\begin{lstlisting}[language=Python, caption=Multimodale Image-Input-Pipeline]
from openai import OpenAI
import base64
from pathlib import Path

class MultimodalPipeline:
    """Verarbeitet Bilder, PDFs und Diagramme via LLM Vision."""

    def __init__(self, model: str = "gpt-4o"):
        self.client = OpenAI()
        self.model = model

    def encode_image(self, image_path: str) -> str:
        """Bild in Base64 kodieren fuer den API-Call."""
        with open(image_path, "rb") as f:
            return base64.b64encode(f.read()).decode("utf-8")

    def analyze_image(self, image_path: str, prompt: str) -> str:
        """Bild mit Prompt analysieren."""
        b64 = self.encode_image(image_path)

        response = self.client.chat.completions.create(
            model=self.model,
            messages=[
                {
                    "role": "user",
                    "content": [
                        {"type": "text", "text": prompt},
                        {
                            "type": "image_url",
                            "image_url": {
                                "url": f"data:image/png;base64,{b64}",
                                "detail": "high",
                            },
                        },
                    ],
                }
            ],
        )
        return response.choices[0].message.content
\end{lstlisting}

\subsection{Vision RAG -- Bilder durchsuchbar machen}

Ein Dokument mit Diagrammen, Screenshots und handschriftlichen Notizen kann nicht direkt in einer Vector-DB gespeichert werden. Die Lösung: \textbf{Vision RAG} -- Bilder werden durch ein multimodales Modell in Text-Beschreibungen übersetzt, die dann gechunkt und embedded werden.

\begin{verbatim}
   Vision-RAG-Pipeline:

   [PDF-Seite] --> [Vision-LLM] --> [Markdown-Beschreibung]
                                        |
                                   [Chunk + Embed]
                                        |
                                   [Vector DB]
                                        |
   [Query] --> [Text-Embedding] --> [Retrieval] --> [LLM-Response]
\end{verbatim}

\begin{lstlisting}[language=Python, caption=Vision-RAG-Pipeline]
from openai import OpenAI


class VisionRAG:
    """RAG-Pipeline fuer bildbasierte Dokumente."""

    def __init__(self, vision_model: str = "gpt-4o",
                 embed_model: str = "text-embedding-3-large"):
        self.vision = MultimodalPipeline(model=vision_model)
        self.embed_client = OpenAI()
        self.embed_model = embed_model
        self.vector_store = {}

    def index_document(self, doc_id: str, image_paths: list[str]):
        """Dokument aus Bildern indexieren."""
        chunks = []
        for i, img_path in enumerate(image_paths):
            # Bild in Markdown uebersetzen
            description = self.vision.analyze_image(
                img_path,
                "Extrahiere den gesamten Text und beschreibe "
                "alle Diagramme, Tabellen und Abbildungen "
                "detailliert. Gib das Ergebnis als Markdown aus."
            )
            # Chunken
            for j, chunk in enumerate(self._chunk_text(description)):
                embedding = self.embed_client.embeddings.create(
                    model=self.embed_model,
                    input=chunk,
                ).data[0].embedding

                chunk_id = f"{doc_id}_p{i}_c{j}"
                self.vector_store[chunk_id] = {
                    "text": chunk,
                    "embedding": embedding,
                    "source": img_path,
                    "page": i,
                }
                chunks.append(chunk_id)

        return chunks

    def search(self, query: str, top_k: int = 3) -> list[dict]:
        query_emb = self.embed_client.embeddings.create(
            model=self.embed_model, input=query
        ).data[0].embedding

        scores = [
            (cid, self._cosine_sim(query_emb, data["embedding"]))
            for cid, data in self.vector_store.items()
        ]
        scores.sort(key=lambda x: -x[1])

        return [
            {"chunk_id": cid, "score": score, **self.vector_store[cid]}
            for cid, score in scores[:top_k]
        ]

    def _chunk_text(self, text: str, size: int = 512) -> list[str]:
        words = text.split()
        return [" ".join(words[i:i+size])
                for i in range(0, len(words), size)]

    def _cosine_sim(self, a: list[float], b: list[float]) -> float:
        import numpy as np
        a, b = np.array(a), np.array(b)
        return float(np.dot(a, b) / (np.linalg.norm(a) * np.linalg.norm(b)))
\end{lstlisting}

\begin{praxishinweis}
Vision RAG lohnt sich ab 50+ Seiten pro Tag. Für kleine Dokumentenmengen reicht es, Bilder direkt per Vision-LLM zu analysieren. Die Embedding-Infrastruktur (Vector-DB, Chunking, Indexing) amortisiert sich erst ab Volumen.
\end{praxishinweis}

\subsection{Diagramme und Charts verarbeiten}

Ein Spezialfall von Vision: technische Diagramme, Architekturskizzen, UML-Diagramme und Charts. Multimodale Modelle sind gut darin, Diagramme zu lesen, aber sie machen systematische Fehler bei Achsenbeschriftungen und Metriken.

\begin{lstlisting}[language=Python, caption=Diagramm-Analyse mit strukturiertem Output]
from openai import OpenAI
from pydantic import BaseModel

class ChartData(BaseModel):
    title: str
    x_axis: str
    y_axis: str
    data_points: list[dict]
    trend: str

def extract_chart_data(chart_path: str) -> ChartData:
    """Extrahiert strukturierte Daten aus einem Chart-Bild."""

    client = OpenAI()
    b64 = base64.b64encode(open(chart_path, "rb").read()).decode()

    completion = client.beta.chat.completions.parse(
        model="gpt-4o",
        messages=[
            {
                "role": "user",
                "content": [
                    {"type": "text",
                     "text": "Extrahiere alle Datenpunkte aus "
                             "diesem Chart als strukturierte Liste."},
                    {
                        "type": "image_url",
                        "image_url": {
                            "url": f"data:image/png;base64,{b64}",
                            "detail": "high",
                        },
                    },
                ],
            }
        ],
        response_format=ChartData,
    )

    return completion.choices[0].message.parsed
\end{lstlisting}

% ===========================================================%
% AUDIO
% ===========================================================%
\section{Audio -- Sprache und Sound}

\subsection{Speech-to-Text (STT)}

Die Umwandlung von Sprache in Text ist die reifste multimodale Technologie. Sie wird praktisch fehlerfrei von GPT-4o-Audio, Whisper v3 und Gemini 3 beherrscht.

\begin{lstlisting}[language=Python, caption=Audio-Transkription mit Whisper]
def transcribe_audio(audio_path: str, language: str = "de") -> dict:
    """Transkribiert eine Audiodatei mit Whisher."""
    from openai import OpenAI

    client = OpenAI()

    with open(audio_path, "rb") as f:
        transcript = client.audio.transcriptions.create(
            model="whisper-1",
            file=f,
            language=language,
            response_format="verbose_json",
            timestamp_granularities=["segment"],
        )

    # Segment-Timestamps fuer Alignment
    segments = [
        {
            "start": seg.start,
            "end": seg.end,
            "text": seg.text.strip(),
        }
        for seg in transcript.segments
    ]

    return {
        "full_text": transcript.text,
        "segments": segments,
        "duration_s": transcript.segments[-1].end,
    }
\end{lstlisting}

\subsection{Multimodales Audio-Processing}

GPT-4o und Gemini 3 verarbeiten Audio direkt, ohne separaten STT-Schritt. Das Modell bekommt die Roh-Audiodaten und kann sowohl Inhalt als auch Tonfall, Emotionen und Sprecherwechsel erkennen.

\begin{lstlisting}[language=Python, caption=Direkte Audio-Verarbeitung (GPT-4o-Audio)]
from openai import OpenAI


class AudioPipeline:
    """Verarbeitet Audio direkt ohne STT-Zwischenschritt."""

    def __init__(self):
        self.client = OpenAI()

    def analyze_audio(self, audio_path: str, prompt: str) -> str:
        """Audio-Inhalt mit Prompt analysieren."""
        with open(audio_path, "rb") as f:
            b64 = base64.b64encode(f.read()).decode()

        response = self.client.chat.completions.create(
            model="gpt-4o-audio-preview",
            messages=[
                {
                    "role": "user",
                    "content": [
                        {"type": "text", "text": prompt},
                        {
                            "type": "input_audio",
                            "input_audio": {
                                "data": b64,
                                "format": "wav",
                            },
                        },
                    ],
                }
            ],
        )
        return response.choices[0].message.content

    def extract_sentiment(self, audio_path: str) -> dict:
        """Stimmung und Sprechermerkmale extrahieren."""
        result = self.analyze_audio(
            audio_path,
            "Analysiere diese Audioaufnahme. "
            "Gib zurueck: Stimmung (positiv/neutral/negativ), "
            "Sprecheranzahl, Hintergrundgeraeusche, "
            "Sprechtempo (woerter/min), Pausen (Anzahl)."
        )
        return result
\end{lstlisting}

\subsection{Text-to-Speech (TTS)}

Die Umkehrung: Sprachsynthese als Output. Wichtig für Voice-Assistenten, Barrierefreiheit und automatisierte Ansagen.

\begin{lstlisting}[language=Python, caption=TTS mit mehreren Stimmen]
def generate_speech(text: str, voice: str = "alloy",
                    output_path: str = "output.mp3") -> str:
    """Erzeugt Sprachausgabe aus Text."""
    from openai import OpenAI

    client = OpenAI()

    response = client.audio.speech.create(
        model="tts-1-hd",
        voice=voice,  # alloy, echo, fable, onyx, nova, shimmer
        input=text,
        speed=1.0,
    )

    response.stream_to_file(output_path)
    return output_path

# Beispiel: Mehrsprachige Ansage
for lang, text, voice in [
    ("de", "Willkommen bei unserem Service.", "nova"),
    ("en", "Welcome to our service.", "alloy"),
    ("fr", "Bienvenue dans notre service.", "shimmer"),
]:
    generate_speech(text, voice=voice,
                    output_path=f"welcome_{lang}.mp3")
\end{lstlisting}

\begin{praxishinweis}
Audio-Native-Verarbeitung (GPT-4o-Audio, Gemini) erhält Emotion und Tonfall, kostet aber mehr Tokens als STT-$\rightarrow$Text-$\rightarrow$LLM. Entscheide nach Use-Case: Sentiment-Analyse braucht natives Audio, Transkription reicht STT.
\end{praxishinweis}

% ===========================================================%
% VIDEO
% ===========================================================%
\section{Video -- Bewegung und Kontext}

\subsection{Frame-Extraktion}

Video-Verarbeitung in LLMs ist noch kein nativer Echtzeit-Stream, sondern erfolgt über \textbf{Frame-Sampling}. Das Modell bekommt eine Auswahl an Frames in regelmässigen Abständen und extrahiert daraus Informationen.

\begin{lstlisting}[language=Python, caption=Video-Frame-Extraktion für LLM-Analyse]
import cv2
import base64
from openai import OpenAI
from pathlib import Path

class VideoAnalyzer:
    """Extrahiert Frames aus Video und analysiert sie multimodal."""

    def __init__(self, model: str = "gpt-4o",
                 frame_interval_s: float = 2.0):
        self.client = OpenAI()
        self.model = model
        self.interval = frame_interval_s

    def extract_frames(self, video_path: str,
                       max_frames: int = 30) -> list[dict]:
        """Extrahiert Frames in regelmaessigen Abstaenden."""
        cap = cv2.VideoCapture(video_path)
        fps = cap.get(cv2.CAP_PROP_FPS)
        frame_interval = int(fps * self.interval)
        frames = []
        frame_count = 0

        while True:
            ret, frame = cap.read()
            if not ret or len(frames) >= max_frames:
                break

            if frame_count % frame_interval == 0:
                _, buffer = cv2.imencode(".jpg", frame, [
                    cv2.IMWRITE_JPEG_QUALITY, 85,
                ])
                b64 = base64.b64encode(buffer).decode()

                timestamp_s = frame_count / fps
                frames.append({
                    "timestamp_s": timestamp_s,
                    "timestamp": f"{int(timestamp_s//60)}m"
                                 f"{int(timestamp_s%60)}s",
                    "base64": b64,
                })

            frame_count += 1

        cap.release()
        return frames

    def analyze_video(self, video_path: str, prompt: str) -> str:
        """Video mit LLM analysieren (Frame-Sampling)."""
        frames = self.extract_frames(video_path)

        # Frames in API-Content umwandeln
        content = [{"type": "text", "text": prompt}]
        for frame in frames:
            content.append({
                "type": "image_url",
                "image_url": {
                    "url": f"data:image/jpeg;base64,{frame['base64']}",
                    "detail": "low",  # low detail fuer viele Frames
                },
            })

        response = self.client.chat.completions.create(
            model=self.model,
            messages=[{"role": "user", "content": content}],
            max_tokens=2048,
        )
        return response.choices[0].message.content
\end{lstlisting}

\subsection{Temporale Reasoning-Strategien}

Video-Analyse erfordert mehr als Einzelframe-Betrachtung. Drei Strategien haben sich bewährt:

\begin{enumerate}[leftmargin=*]
    \item \textbf{Frame-Stacking}: Alle Frames auf einmal senden. Gut für kurze Clips ($<$ 30s), aber teuer (30 Frames $\times$ 85 Tokens Basis + Tiling).
    \item \textbf{Zusammenfassungs-Pipeline}: Frame-Gruppen separat analysieren, Ergebnisse aggregieren. Gut für lange Videos.
    \item \textbf{Keyframe-Detektion}: Nur Frames mit signifikanten Änderungen extrahieren (Szene-Wechsel, Bewegung).
\end{enumerate}

\begin{lstlisting}[language=Python, caption=Hierarchische Video-Analyse]
class HierarchicalVideoAnalyzer(VideoAnalyzer):
    """Zusammenfassungs-Pipeline fuer lange Videos."""

    def analyze_long_video(self, video_path: str) -> dict:
        # Phase 1: Segmente analysieren (alle 10s ein Frame)
        self.interval = 10.0
        frames = self.extract_frames(video_path, max_frames=50)

        # Phase 2: Frames in Gruppen zu je 5 teilen
        segments = []
        for i in range(0, len(frames), 5):
            group = frames[i:i+5]
            segment_content = [
                {"type": "text",
                 "text": "Fasse zusammen, was in diesen "
                         "Videosegmenten passiert."},
            ]
            for f in group:
                segment_content.append({
                    "type": "image_url",
                    "image_url": {
                        "url": f"data:image/jpeg;base64,{f['base64']}",
                        "detail": "low",
                    },
                })

            resp = self.client.chat.completions.create(
                model=self.model,
                messages=[{"role": "user",
                          "content": segment_content}],
                max_tokens=512,
            )
            segments.append({
                "segment": i // 5,
                "timestamp": group[0]["timestamp"],
                "summary": resp.choices[0].message.content,
            })

        # Phase 3: Gesamt-Zusammenfassung
        all_summaries = "\n".join(
            f"[{s['timestamp']}] {s['summary']}"
            for s in segments
        )
        final = self.client.chat.completions.create(
            model=self.model,
            messages=[
                {"role": "user",
                 "content": f"Erstelle eine zusammenhaengende "
                            f"Zusammenfassung dieses Videos:\n\n"
                            f"{all_summaries}"}
            ],
        )
        return {
            "segments": segments,
            "summary": final.choices[0].message.content,
            "total_segments": len(segments),
        }
\end{lstlisting}

\begin{praxishinweis}
Frame-Sampling ist ein Kosten-Kompromiss. 1 Frame/s für Szenen-Wechsel, 1 Frame/10s für statische Präsentationen. Und immer low-detail für Video-Frames -- high-detail kostet 4x und bringt bei Bewegung kaum Mehrwert.
\end{praxishinweis}

% ===========================================================%
% BEST PRACTICES
% ===========================================================%
\section{Best Practices}

\begin{itemize}[leftmargin=*]
    \item \textbf{Detail-Level wählen}: \texttt{high} detail kostet 4x mehr Tokens als \texttt{low}. Nutze low für Screenshots, high für feine Diagramme.
    \item \textbf{Bilder vorverarbeiten}: Zuschneiden auf relevante Regionen, Rauschen reduzieren, Kontrast optimieren. Ein bearbeitetes Bild kostet gleich viele Tokens, liefert bessere Ergebnisse.
    \item \textbf{Frame-Intervalle anpassen}: Für Szenen-Wechsel 1 Frame/s, für statische Präsentationen 1 Frame/10s reicht.
    \item \textbf{Audio vor Segmentierung}: Lange Audiodateien ($>$ 30 Min) vor der Verarbeitung in sinnvolle Segmente teilen (Stille als Trennmarker).
    \item \textbf{Caching für Bilder}: Gleiche Bilder (Logos, Icons, Header) werden immer gleich interpretiert -- Cache den Vision-Output.
\end{itemize}

% ===========================================================%
% ANTI-PATTERNS
% ===========================================================%
\section{Anti-Patterns}

\begin{itemize}[leftmargin=*]
    \item \textbf{Jedes Bild in High Resolution}: Ein 4K-Bild kostet 40+ Tiles = $>$ 7.000 Tokens. Meist reicht Low-Res.
    \item \textbf{Video-Frame-Overload}: 60 Frames an GPT-4o senden kostet 6.800+ Tokens nur für das Bildmaterial. Nutze hierarchische Zusammenfassung.
    \item \textbf{Bilderschrift lesen ohne OCR-Check}: LLMs lesen Text in Bildern gut, aber nicht perfekt. Bei 100\% Genauigkeitsanforderung: Tesseract-OCR vor LLM-Analyse schalten.
    \item \textbf{Audio ohne Timestamps}: Transkription ohne Segment-Timestamps macht die Zuordnung zu Video-Frames unmöglich. Immer \texttt{timestamp\_granularities=["segment"]} setzen.
\end{itemize}

% ===========================================================%
% SICHERHEIT
% ===========================================================%
\section{Sicherheit multimodaler Systeme}

Multimodale Eingaben eröffnen neue Angriffsvektoren:

\begin{enumerate}[leftmargin=*]
    \item \textbf{Prompt Injection via Bild}: Ein Angreifer bettet unsichtbaren Text in ein Bild ein (z.B. weisse Schrift auf weissem Grund), der vom Vision-Encoder gelesen wird.
    \item \textbf{OCR-Exploits}: Text in Bildern wird nicht durch Input-Filter erkannt. Ein klassischer Prompt-Injection-Filter greift nicht, wenn der Injection-Text im Bild steht.
    \item \textbf{Adversarial Patches}: Kleine Bildregionen können das Modell gezielt manipulieren (falsche Klassifikation, übersehene Objekte).
    \item \textbf{Acoustic Attacks}: Ultraschalltöne, die für Menschen unhörbar sind, können STT-Systeme manipulieren.
\end{enumerate}

\begin{lstlisting}[language=Python, caption=Multimodaler Input-Guard]
class MultimodalGuard:
    """Sicherheitsfilter fuer multimodale Eingaben."""

    def __init__(self):
        self.text_filter = InputGuard()  # aus Kapitel 19

    def check_image(self, image_b64: str) -> bool:
        """Prueft Bild auf Injection-Versuche."""
        import pytesseract
        from PIL import Image
        import io

        image = Image.open(io.BytesIO(base64.b64decode(image_b64)))

        # OCR-extrahierter Text
        ocr_text = pytesseract.image_to_string(image)

        # Injection-Check im OCR-Text
        suspicious = self.text_filter.check(ocr_text)
        if suspicious:
            logger.warning({
                "event": "image_injection_detected",
                "ocr_text": ocr_text[:200],
                "reason": suspicious,
            })
            return False

        return True

    def check_audio(self, audio_path: str) -> bool:
        """Prueft Audio auf ungewoehnliche Frequenzen."""
        import librosa

        y, sr = librosa.load(audio_path, sr=None)
        # Pruefung auf Ultraschall (Frequenzen > 20 kHz)
        # Vereinfachte Implementierung
        return True
\end{lstlisting}

\begin{praxishinweis}
Multimodale Sicherheit ist kein nachträglicher Einbau. OCR-Injection-Check und Audio-Frequenz-Filter gehören in die Basis-Pipeline, nicht ins Monitoring-Dashboard. Baut den Input Guard bevor der erste User den Endpoint testet.
\end{praxishinweis}

% ===========================================================%
% ZUSAMMENFASSUNG
% ===========================================================%
\begin{zusammenfassung}
\begin{itemize}[leftmargin=*]
    \item Multimodale Modelle (GPT-4o, Claude 4 Vision, Gemini 3) verarbeiten Text, Bilder, Audio und Video in einem Modell.
    \item Bilder werden in Tokens umgerechnet -- ein 1024$\times$1024-Bild kostet ca. 2.000 Tokens. Das Detail-Level (high/low) beeinflusst Kosten und Qualität.
    \item Vision RAG übersetzt Bilder in Text-Beschreibungen, die dann gechunkt und embedded werden.
    \item Video-Analyse erfolgt über Frame-Sampling. Für lange Videos ist eine hierarchische Zusammenfassung (Segment $\rightarrow$ Gesamt) nötig.
    \item Multimodale Systeme erfordern neue Sicherheitsmassnahmen: OCR-Injection-Check, adversariale Patches, Audio-Frequenz-Analyse.
\end{itemize}
\end{zusammenfassung}

\begin{merke}
Multimodalität ist kein Feature -- es ist ein neues Paradigma. Die Pipelines, Sicherheitsmodelle und Cost-Strategien, die für Text LLMs entwickelt wurden, lassen sich nicht eins-zu-eins übertragen. Jede Modalität hat ihre eigenen Kostenstrukturen, Fehlerbilder und Angriffsvektoren.
\end{merke}

% ===========================================================%
% PRAXISPROJEKT
% ===========================================================%
\section{Praxisprojekt -- Multimodaler Dokumenten-Assistent}

\textbf{Ziel:} Baue einen Assistenten, der gescannte PDF-Dokumente (Rechnungen, Verträge, Formulare) verarbeitet, strukturiert extrahiert und per RAG durchsuchbar macht.

\textbf{Aufgaben:}
\begin{enumerate}[leftmargin=*]
    \item PDF-Seiten als Bilder extrahieren (PyMuPDF)
    \item Mit Vision-LLM in Markdown übersetzen
    \item Markdown chunken + embedden + in Vector DB speichern
    \item Extraktion: \texttt{response\_format} mit Pydantic-Schema für strukturierte Daten (Rechnungsnummer, Datum, Betrag, Adresse)
    \item RAG-Query-Interface mit Quellenangabe pro Seite
\end{enumerate}

\textbf{Erfolgskriterium:} Eine Rechnung wird in 5 Sekunden indexiert und ist per Natürlichsprach-Query in 100\,\% der Fälle korrekt auffindbar.

% ===========================================================%
% RESSOURCEN
% ===========================================================%
\section{Ressourcen}

\begin{itemize}[leftmargin=*]
    \item \textbf{OpenAI Vision Guide}: \href{https://platform.openai.com/docs/guides/vision}{platform.openai.com/docs/guides/vision}
    \item \textbf{OpenAI Audio Guide}: \href{https://platform.openai.com/docs/guides/audio}{platform.openai.com/docs/guides/audio}
    \item \textbf{Claude Vision}: \href{https://docs.anthropic.com/en/docs/build-with-claude/vision}{docs.anthropic.com/en/docs/build-with-claude/vision}
    \item \textbf{Gemini Vision/Video}: \href{https://ai.google.dev/gemini-api/docs/vision}{ai.google.dev/gemini-api/docs/vision}
    \item \textbf{OpenCV Tutorials}: \href{https://docs.opencv.org/}{docs.opencv.org}
    \item \textbf{Tesseract OCR}: \href{https://github.com/tesseract-ocr/tesseract}{github.com/tesseract-ocr/tesseract}
    \item \textbf{PyMuPDF}: \href{https://pymupdf.readthedocs.io/}{pymupdf.readthedocs.io}
\end{itemize}

\textbf{Nächster Schritt:} Im nächsten Teil des Buches tauchst du in den Production Layer ein -- beginnend mit Inference Optimization und Serving für skalierbare LLM-Anwendungen.
